\scanpage{44}
Tout ss-gpe lisse \inserted{\unsure{vert$^{\mathrm t}$|virt$^{\mathrm t}$|rel.} connexe} de \deleted{rang} dim rel. 1 dans $G$ est-il de la forme $C(\mathfrak h)$ -- i.e. est-il déterminé par $\operatorname{Lie}H$ ? Oui \inserted{s'il est \uline{comm} et} si les car résiduelles sont $\ne2$, car il est contenu alors dans $C(\mathfrak h)$ qui a les $=$ vertus, donc lui est égal. On se demande si la condition de commutativité est automatiquement satisfaite ?

\deleted{\unsure{D'où|D'ail.}} \deleted{Question infinitésimale} On sait que c'est le cas sur un corps -- tt grpe lisse \inserted{conn.} de dim 1 est comm. -- donc la question est une question de déformations de ss-groupes, liée à un $H^2$ de Hochschild --

Considérons le normalisateur de $H$ -- qui est aussi celui de $\mathfrak h$, puisque la corr. entre $H$ et $\mathfrak h$ est $1$-$1$. J'ai envie de prouver qu'il est (du moins en car rés. $\ne2$) lisse de dim relative 1 -- mais c'est faux $=$ sur un corps, quand $H$ est unipotent, son normalisateur devient un Borel ! C'est ici qu'il faut faire \unsure{intervenir|interven}
% archive sheet number: 43
